% Options for packages loaded elsewhere
% Options for packages loaded elsewhere
\PassOptionsToPackage{unicode,linktoc=all}{hyperref}
\PassOptionsToPackage{hyphens}{url}
\PassOptionsToPackage{dvipsnames,svgnames,x11names}{xcolor}
%
\documentclass[
  11pt,
  a4paper,
  twoside]{scrbook}
\usepackage{xcolor}
\usepackage[inner=3cm,outer=4cm,top=2.5cm,bottom=2.5cm,left=3cm,right=1.5cm,headsep=22pt,headheight=11pt,footskip=33pt,ignorehead,ignorefoot,heightrounded]{geometry}
\usepackage{amsmath,amssymb}
\setcounter{secnumdepth}{2}
\usepackage{iftex}
\ifPDFTeX
  \usepackage[T1]{fontenc}
  \usepackage[utf8]{inputenc}
  \usepackage{textcomp} % provide euro and other symbols
\else % if luatex or xetex
  \usepackage{unicode-math} % this also loads fontspec
  \defaultfontfeatures{Scale=MatchLowercase}
  \defaultfontfeatures[\rmfamily]{Ligatures=TeX,Scale=1}
\fi
\usepackage{lmodern}
\ifPDFTeX\else
  % xetex/luatex font selection
  \setmainfont[]{Latin Modern Roman}
\fi
% Use upquote if available, for straight quotes in verbatim environments
\IfFileExists{upquote.sty}{\usepackage{upquote}}{}
\IfFileExists{microtype.sty}{% use microtype if available
  \usepackage[]{microtype}
  \UseMicrotypeSet[protrusion]{basicmath} % disable protrusion for tt fonts
}{}
\makeatletter
\@ifundefined{KOMAClassName}{% if non-KOMA class
  \IfFileExists{parskip.sty}{%
    \usepackage{parskip}
  }{% else
    \setlength{\parindent}{0pt}
    \setlength{\parskip}{6pt plus 2pt minus 1pt}}
}{% if KOMA class
  \KOMAoptions{parskip=half}}
\makeatother
% Make \paragraph and \subparagraph free-standing
\makeatletter
\ifx\paragraph\undefined\else
  \let\oldparagraph\paragraph
  \renewcommand{\paragraph}{
    \@ifstar
      \xxxParagraphStar
      \xxxParagraphNoStar
  }
  \newcommand{\xxxParagraphStar}[1]{\oldparagraph*{#1}\mbox{}}
  \newcommand{\xxxParagraphNoStar}[1]{\oldparagraph{#1}\mbox{}}
\fi
\ifx\subparagraph\undefined\else
  \let\oldsubparagraph\subparagraph
  \renewcommand{\subparagraph}{
    \@ifstar
      \xxxSubParagraphStar
      \xxxSubParagraphNoStar
  }
  \newcommand{\xxxSubParagraphStar}[1]{\oldsubparagraph*{#1}\mbox{}}
  \newcommand{\xxxSubParagraphNoStar}[1]{\oldsubparagraph{#1}\mbox{}}
\fi
\makeatother

\usepackage{color}
\usepackage{fancyvrb}
\newcommand{\VerbBar}{|}
\newcommand{\VERB}{\Verb[commandchars=\\\{\}]}
\DefineVerbatimEnvironment{Highlighting}{Verbatim}{commandchars=\\\{\}}
% Add ',fontsize=\small' for more characters per line
\newenvironment{Shaded}{}{}
\newcommand{\AlertTok}[1]{\textcolor[rgb]{1.00,0.00,0.00}{\textbf{#1}}}
\newcommand{\AnnotationTok}[1]{\textcolor[rgb]{0.38,0.63,0.69}{\textbf{\textit{#1}}}}
\newcommand{\AttributeTok}[1]{\textcolor[rgb]{0.49,0.56,0.16}{#1}}
\newcommand{\BaseNTok}[1]{\textcolor[rgb]{0.25,0.63,0.44}{#1}}
\newcommand{\BuiltInTok}[1]{\textcolor[rgb]{0.00,0.50,0.00}{#1}}
\newcommand{\CharTok}[1]{\textcolor[rgb]{0.25,0.44,0.63}{#1}}
\newcommand{\CommentTok}[1]{\textcolor[rgb]{0.38,0.63,0.69}{\textit{#1}}}
\newcommand{\CommentVarTok}[1]{\textcolor[rgb]{0.38,0.63,0.69}{\textbf{\textit{#1}}}}
\newcommand{\ConstantTok}[1]{\textcolor[rgb]{0.53,0.00,0.00}{#1}}
\newcommand{\ControlFlowTok}[1]{\textcolor[rgb]{0.00,0.44,0.13}{\textbf{#1}}}
\newcommand{\DataTypeTok}[1]{\textcolor[rgb]{0.56,0.13,0.00}{#1}}
\newcommand{\DecValTok}[1]{\textcolor[rgb]{0.25,0.63,0.44}{#1}}
\newcommand{\DocumentationTok}[1]{\textcolor[rgb]{0.73,0.13,0.13}{\textit{#1}}}
\newcommand{\ErrorTok}[1]{\textcolor[rgb]{1.00,0.00,0.00}{\textbf{#1}}}
\newcommand{\ExtensionTok}[1]{#1}
\newcommand{\FloatTok}[1]{\textcolor[rgb]{0.25,0.63,0.44}{#1}}
\newcommand{\FunctionTok}[1]{\textcolor[rgb]{0.02,0.16,0.49}{#1}}
\newcommand{\ImportTok}[1]{\textcolor[rgb]{0.00,0.50,0.00}{\textbf{#1}}}
\newcommand{\InformationTok}[1]{\textcolor[rgb]{0.38,0.63,0.69}{\textbf{\textit{#1}}}}
\newcommand{\KeywordTok}[1]{\textcolor[rgb]{0.00,0.44,0.13}{\textbf{#1}}}
\newcommand{\NormalTok}[1]{#1}
\newcommand{\OperatorTok}[1]{\textcolor[rgb]{0.40,0.40,0.40}{#1}}
\newcommand{\OtherTok}[1]{\textcolor[rgb]{0.00,0.44,0.13}{#1}}
\newcommand{\PreprocessorTok}[1]{\textcolor[rgb]{0.74,0.48,0.00}{#1}}
\newcommand{\RegionMarkerTok}[1]{#1}
\newcommand{\SpecialCharTok}[1]{\textcolor[rgb]{0.25,0.44,0.63}{#1}}
\newcommand{\SpecialStringTok}[1]{\textcolor[rgb]{0.73,0.40,0.53}{#1}}
\newcommand{\StringTok}[1]{\textcolor[rgb]{0.25,0.44,0.63}{#1}}
\newcommand{\VariableTok}[1]{\textcolor[rgb]{0.10,0.09,0.49}{#1}}
\newcommand{\VerbatimStringTok}[1]{\textcolor[rgb]{0.25,0.44,0.63}{#1}}
\newcommand{\WarningTok}[1]{\textcolor[rgb]{0.38,0.63,0.69}{\textbf{\textit{#1}}}}

\usepackage{longtable,booktabs,array}
\newcounter{none} % for unnumbered tables
\usepackage{calc} % for calculating minipage widths
% Correct order of tables after \paragraph or \subparagraph
\usepackage{etoolbox}
\makeatletter
\patchcmd\longtable{\par}{\if@noskipsec\mbox{}\fi\par}{}{}
\makeatother
% Allow footnotes in longtable head/foot
\IfFileExists{footnotehyper.sty}{\usepackage{footnotehyper}}{\usepackage{footnote}}
\makesavenoteenv{longtable}
\usepackage{graphicx}
\makeatletter
\newsavebox\pandoc@box
\newcommand*\pandocbounded[1]{% scales image to fit in text height/width
  \sbox\pandoc@box{#1}%
  \Gscale@div\@tempa{\textheight}{\dimexpr\ht\pandoc@box+\dp\pandoc@box\relax}%
  \Gscale@div\@tempb{\linewidth}{\wd\pandoc@box}%
  \ifdim\@tempb\p@<\@tempa\p@\let\@tempa\@tempb\fi% select the smaller of both
  \ifdim\@tempa\p@<\p@\scalebox{\@tempa}{\usebox\pandoc@box}%
  \else\usebox{\pandoc@box}%
  \fi%
}
% Set default figure placement to htbp
\def\fps@figure{htbp}
\makeatother


% definitions for citeproc citations
\NewDocumentCommand\citeproctext{}{}
\NewDocumentCommand\citeproc{mm}{%
  \begingroup\def\citeproctext{#2}\cite{#1}\endgroup}
\makeatletter
 % allow citations to break across lines
 \let\@cite@ofmt\@firstofone
 % avoid brackets around text for \cite:
 \def\@biblabel#1{}
 \def\@cite#1#2{{#1\if@tempswa , #2\fi}}
\makeatother
\newlength{\cslhangindent}
\setlength{\cslhangindent}{1.5em}
\newlength{\csllabelwidth}
\setlength{\csllabelwidth}{3em}
\newenvironment{CSLReferences}[2] % #1 hanging-indent, #2 entry-spacing
 {\begin{list}{}{%
  \setlength{\itemindent}{0pt}
  \setlength{\leftmargin}{0pt}
  \setlength{\parsep}{0pt}
  % turn on hanging indent if param 1 is 1
  \ifodd #1
   \setlength{\leftmargin}{\cslhangindent}
   \setlength{\itemindent}{-1\cslhangindent}
  \fi
  % set entry spacing
  \setlength{\itemsep}{#2\baselineskip}}}
 {\end{list}}
\usepackage{calc}
\newcommand{\CSLBlock}[1]{\hfill\break\parbox[t]{\linewidth}{\strut\ignorespaces#1\strut}}
\newcommand{\CSLLeftMargin}[1]{\parbox[t]{\csllabelwidth}{\strut#1\strut}}
\newcommand{\CSLRightInline}[1]{\parbox[t]{\linewidth - \csllabelwidth}{\strut#1\strut}}
\newcommand{\CSLIndent}[1]{\hspace{\cslhangindent}#1}



\setlength{\emergencystretch}{3em} % prevent overfull lines

\providecommand{\tightlist}{%
  \setlength{\itemsep}{0pt}\setlength{\parskip}{0pt}}



 


% ---------- PROFESSOR'S HEADER (KEEP THIS) ----------
\addtokomafont{disposition}{\rmfamily}
\KOMAoptions{chapterprefix=true,appendixprefix=true}
\KOMAoptions{headings=small}

\setkomafont{pageheadfoot}{\normalfont\normalcolor\footnotesize}
\setkomafont{pagenumber}{\normalfont\normalcolor\footnotesize}

\usepackage{scrlayer-scrpage}
\lefoot[]{}
\cefoot[\pagemark]{}
\refoot[]{}
\lofoot[]{}
\cofoot[\pagemark]{}
\rofoot[]{}
\lehead[]{\pagemark}
\cehead[]{}
\rehead[]{\leftmark}
\lohead[]{\rightmark}
\cohead[]{}
\rohead[]{\pagemark}


% ---------- OPTIONAL IMPROVEMENTS ----------
% No blank pages before chapters in pdf
\KOMAoptions{open=any, cleardoublepage=empty}

% Better chapter spacing
\RedeclareSectionCommand[
  beforeskip=-1ex plus -1ex minus -1ex,
  afterskip=2ex plus 1ex minus 1ex
]{chapter}
\makeatletter
\@ifpackageloaded{tcolorbox}{}{\usepackage[skins,breakable]{tcolorbox}}
\@ifpackageloaded{fontawesome5}{}{\usepackage{fontawesome5}}
\definecolor{quarto-callout-color}{HTML}{909090}
\definecolor{quarto-callout-note-color}{HTML}{0758E5}
\definecolor{quarto-callout-important-color}{HTML}{CC1914}
\definecolor{quarto-callout-warning-color}{HTML}{EB9113}
\definecolor{quarto-callout-tip-color}{HTML}{00A047}
\definecolor{quarto-callout-caution-color}{HTML}{FC5300}
\definecolor{quarto-callout-color-frame}{HTML}{acacac}
\definecolor{quarto-callout-note-color-frame}{HTML}{4582ec}
\definecolor{quarto-callout-important-color-frame}{HTML}{d9534f}
\definecolor{quarto-callout-warning-color-frame}{HTML}{f0ad4e}
\definecolor{quarto-callout-tip-color-frame}{HTML}{02b875}
\definecolor{quarto-callout-caution-color-frame}{HTML}{fd7e14}
\makeatother
\makeatletter
\@ifpackageloaded{bookmark}{}{\usepackage{bookmark}}
\makeatother
\makeatletter
\@ifpackageloaded{caption}{}{\usepackage{caption}}
\AtBeginDocument{%
\ifdefined\contentsname
  \renewcommand*\contentsname{Table of contents}
\else
  \newcommand\contentsname{Table of contents}
\fi
\ifdefined\listfigurename
  \renewcommand*\listfigurename{List of Figures}
\else
  \newcommand\listfigurename{List of Figures}
\fi
\ifdefined\listtablename
  \renewcommand*\listtablename{List of Tables}
\else
  \newcommand\listtablename{List of Tables}
\fi
\ifdefined\figurename
  \renewcommand*\figurename{Figure}
\else
  \newcommand\figurename{Figure}
\fi
\ifdefined\tablename
  \renewcommand*\tablename{Table}
\else
  \newcommand\tablename{Table}
\fi
}
\@ifpackageloaded{float}{}{\usepackage{float}}
\floatstyle{ruled}
\@ifundefined{c@chapter}{\newfloat{codelisting}{h}{lop}}{\newfloat{codelisting}{h}{lop}[chapter]}
\floatname{codelisting}{Listing}
\newcommand*\listoflistings{\listof{codelisting}{List of Listings}}
\makeatother
\makeatletter
\makeatother
\makeatletter
\@ifpackageloaded{caption}{}{\usepackage{caption}}
\@ifpackageloaded{subcaption}{}{\usepackage{subcaption}}
\makeatother
\usepackage{bookmark}
\IfFileExists{xurl.sty}{\usepackage{xurl}}{} % add URL line breaks if available
\urlstyle{same}
% fallback for those not using the hyperref driver hyperxmp:
\makeatletter
\@ifundefined{xmpquote}{\newcommand{\xmpquote}[1]{#1}}{}
\makeatother
\hypersetup{
  pdftitle={Evaluation of LABOR Synthesizer for Educational Purposes in Electronics},
  pdfauthor={Punnavarjula Sai Harshita},
  colorlinks=true,
  linkcolor={blue},
  filecolor={Maroon},
  citecolor={Blue},
  urlcolor={Blue},
  pdfcreator={LaTeX via pandoc}}


\title{Evaluation of LABOR Synthesizer for Educational Purposes in
Electronics}
\author{Punnavarjula Sai Harshita}
\date{October 2026}
\begin{document}
\frontmatter
\maketitle

\begin{titlepage}
\begin{center}

% ------------------------------------------------
% LOGO
% ------------------------------------------------
\vspace*{-1cm}
\includegraphics[width=5.5cm]{images/hsb-logo.png}

\vspace{0.8cm}

{\large Faculty of Electrical Engineering and Computer Science}

\vspace{1cm}

% ------------------------------------------------
% TITLE
% ------------------------------------------------
{\bfseries\Large
Evaluation of LABOR Synthesizer for Educational Purposes in Electronics
}\\[1cm]




% ------------------------------------------------
% AUTHORS
% ------------------------------------------------
{\large by}\\[10pt]

{\Large Punnavarjula Sai Harshita }\\
{\large Matr.-No.: }\\[16pt]

\vspace{1cm}

% ------------------------------------------------
% SUBMISSION INFO
% ------------------------------------------------
{\large A Thesis presented to the City University Bremen \\
in fulfillment of the thesis requirement for the degree of \\
Master in Electronics Engineering}

\vspace{1cm}

% ------------------------------------------------
% EXAMINERS
% ------------------------------------------------
\begin{tabular}{rl}
1. Examiner: & Prof. Dr.-Ing. M. Meiners \\
2. Examiner: & Prof. Dr.-Ing. M. Goldenbaum \\
\end{tabular}

\vspace{1cm}

% ------------------------------------------------
% DATES
% ------------------------------------------------
{\large Preparation Period: 10.2026 -- 05.2027}

\end{center}
\end{titlepage}

% \mainmatter
% \clearpage

\renewcommand*\contentsname{Table of contents}
{
\setcounter{tocdepth}{2}
\tableofcontents
}
\listoffigures
\listoftables

\mainmatter
\bookmarksetup{startatroot}

\chapter*{Preface}\label{preface}
\addcontentsline{toc}{chapter}{Preface}

\markboth{Preface}{Preface}

This thesis template is intended for honours, masters or PhD students at
the Australian National University (ANU) who wish to write their thesis
using the \href{https://quarto.org/}{Quarto} document format. It is
highly recommended for students who code using Python, R or Julia and
have many computational or analysis results in their thesis.

\begin{tcolorbox}[enhanced jigsaw, arc=.35mm, bottomrule=.15mm, bottomtitle=1mm, breakable, colback=white, colbacktitle=quarto-callout-note-color!10!white, colframe=quarto-callout-note-color-frame, coltitle=black, left=2mm, leftrule=.75mm, opacityback=0, opacitybacktitle=0.6, rightrule=.15mm, title=\textcolor{quarto-callout-note-color}{\faInfo}\hspace{0.5em}{Note}, titlerule=0mm, toprule=.15mm, toptitle=1mm]

This thesis template is available on GitHub at
\href{https://github.com/anuopensci/quarto-anu-thesis}{github.com/anuopensci/quarto-anu-thesis}.

\end{tcolorbox}

\section*{Benefits}\label{benefits}
\addcontentsline{toc}{section}{Benefits}

\markright{Benefits}

The benefits of using Quarto document include:

\begin{itemize}
\tightlist
\item
  It allows you to write your thesis in a simple markup language called
  \href{https://www.markdownguide.org/}{Markdown}. This means that you
  can focus on writing your thesis without having to worry about
  formatting.
\item
  The document can be output to a variety of formats including PDF,
  HTML, Word, and LaTeX.
\item
  Code can be easily embedded in the document and executed. This means
  that you can include the results of your analysis in your thesis
  without having to manually copy and paste them. This is a good
  reproducible and scientific practice.
\item
  You can easily integrate with aspects of GitHub (edit, reporting an
  issue, etc).
\end{itemize}

The above outlined benefits can also be considered as best practice.
Version controlling and collaborative writing (via Git and GitHub) are
important in managing multiple versions of your thesis and in
collaborating with your supervisory team. Embedding code in your thesis
is a good practice in reproducible research. Making your thesis in HTML
format can allow for interactive web elements to be embedded while PDF
format can be for general distribution and printing.

\section*{Getting started}\label{getting-started}
\addcontentsline{toc}{section}{Getting started}

\markright{Getting started}

There are several systems that you are expected to know to use this
template. These include:

\begin{itemize}
\tightlist
\item
  Markdown syntax for writing
\item
  Quarto or R Markdown syntax (note that these works for Python or Julia
  too) for embedding code
\item
  (Optional) Git and GitHub for hosting
\end{itemize}

\section*{Frequently asked questions}\label{frequently-asked-questions}
\addcontentsline{toc}{section}{Frequently asked questions}

\markright{Frequently asked questions}

\subsection*{What about Overleaf?}\label{what-about-overleaf}
\addcontentsline{toc}{subsection}{What about Overleaf?}

ANU has a professional account for Overleaf, which is great for those
that use LaTeX regularly. Unfortunately, there is no equivalent system
with track changes in Quarto. You can output the tex file from Quarto
document and use this in Overleaf. The changes made in this tex document
however has to be manually transferred back to the Quarto document. If
your main output is mainly mathematical and you have little to no code
outputs, Overleaf is probably a better choice.

\bookmarksetup{startatroot}

\chapter*{Declaration of Authorship}\label{declaration-of-authorship}

\markboth{Declaration of Authorship}{Declaration of Authorship}

\url{https://www.hs-bremen.de/assets/hsb/de/Dokumente/ZLL/MMCC/declaration_indep_work_Nov_2024_Form.pdf}

I hereby certify that I have written this thesis independently and have
not used any sources or aids other than those specified. Sources and
aids other than those indicated. The passages in the thesis that have
been taken from other works in terms of wording or meaning are from
other works are identified as such.

This statement also applies to the graphics, sketches and illustrations
contained in the thesis. I have not already submitted the work in the
same or a similar form, including extracts, as part of an examination or
coursework.

I hereby certify that the electronic version of the dissertation
submitted is a complete copy of the printed version. and agree to an
electronic check of the thesis using plagiarism software.

Please tick the box: One of the two options must be chosen in
consultation between the assessor and the assessee.

\begin{itemize}
\item[$\square$]
  Option 1: Use of AI-based tools without mandatory labelling I confirm
  that we have not used any AI-based tools whose use has not been agreed
  in writing with the verifier. I understand that the use of text or
  other content and products generated by AI-based tools is not a
  guarantee of their quality. I accept responsibility for the work
  submitted and the content contained therein. I also assure you that my
  creative influence predominates in this work.
\item[$\square$]
  Option 2: Use of AI-based tools with mandatory labelling I confirm
  that we have not used any AI-based tools whose use has not been agreed
  in writing with the verifier. I understand that the use of text or
  other content and products generated by AI-based tools is not a
  guarantee of their quality. I accept responsibility for the work
  submitted and the content contained therein. I also assure you that my
  creative influence predominates in this work. All verbatim or
  analogous adaptations and quotations, and all sections designed,
  written and/or edited using AI-based tools, will be marked and checked
  by me. The form of marking will be agreed between the examiner and the
  examinee.
\end{itemize}

We are aware that failure to comply with the above may have consequences
in terms of examination law, and in particular may result in the
examination being marked ``insufficient'' or the coursework being marked
``failed'' and, in the case of multiple or serious attempts to cheat,
may result in my being de-registered.

First name and surname:

Matriculation number:

Date:

Signature:

\bookmarksetup{startatroot}

\chapter*{Abstract}\label{abstract}
\addcontentsline{toc}{chapter}{Abstract}

\markboth{Abstract}{Abstract}

\bookmarksetup{startatroot}

\chapter*{Acknowledgements}\label{acknowledgements}
\addcontentsline{toc}{chapter}{Acknowledgements}

\markboth{Acknowledgements}{Acknowledgements}

I would like to express my sincere gratitude to my dog, Chuckles, for
eating my research notes multiple times. If it wasn't for you, I would
have finished this thesis earlier.

\bookmarksetup{startatroot}

\chapter{Introduction}\label{sec-intro}

Synthesizer is an electronic musical instrument that generates audio
signals which can be converted to sound. Synthesizers create a wide
variety of sounds by manipulating electrical signals through various
synthesis techniques. Analog synthesizers, in particular, use continuous
voltage signals to produce sound. The design and implementation of
analog synthesizer circuits require a deep understanding of electronic
components, signal processing, and circuit design principles.

Instead of adding several components to a breadboard and connecting them
with wires, the LABOR workstation provides a more structured and
efficient environment for circuit prototyping. It allows for the
integration of multiple circuit modules which includes a full-size
breadboard, a dual power supply with short-circuit protection, an
oscilloscope, and built-in signal and control sources and an interfacing
section, enabling the creation of complex synthesizer architectures
without the need for excessive wiring. This makes it an ideal platform
for educational purposes, allowing students to focus on understanding
the underlying principles of analog synthesis and known as ``All in one
workstation''.

This thesis presents the design and implementation of electronic circuit
architectures for analog music synthesis using the LABOR workstation as
a practical development and prototyping platform. The work integrates
hardware circuit design with software components to create a functional
electronic music synthesis system suitable for educational purposes
particularly for observing, measuring and evaluating the generated
signals {[}1{]}.

The work will be inspired by Moritz Klein's synthesizer-design tutorials
and approach starts with fundamental circuit concepts and progressively
develops them into functional synthesizer building blocks {[}2{]}.

The goal is not only to reproduce existing synthesizer circuits but
rather their underlying design ideas will be used as a starting point
for developing, testing and documenting an independently designed set of
circuits on the LABOR workstation.

\section{Overview}\label{overview}

This thesis will explore three circuit areas that can be combined to
create a small experimental electronic music system. The first part will
focus on converting a sawtooth waveform into a triangular waveform, with
particular attention to how the circuit design and component values
influence the output. The second part will involve developing a simple
hardware-based sequencer for generating rhythmic patterns using clock
signals, shift registers, and sequence controls. The final part will
investigate the generation of different types of noise, particularly
white, pink, and blue noise, and how their frequency characteristics can
be shaped through filtering. After developing and testing the three
circuit sections individually.@Meyer\_1976 {[}3{]} {[}4{]} {[}5{]}
{[}6{]}

The main purpose of the thesis is not only to build the circuits, but
also gain a practical understanding of how they work and how different
components and design choices affect their behaviour. The circuits will
be tested in the laboratory, and the measured results will be compared
with the expected theoretical behaviour. Software tools will also be
used when they are useful for visualising or analysing the signals.

This approach will be expected to help students connect theoretical
principles with observable sound phenomena and measurable electrical
signals. In each exercise, students predict the result, measure it, and
compare it with theory. The built-in tools of the LABOR workstation
supports the experimental process by minimizing potential wiring errors
and enabling students to carry out the exercises with little
supervision. Designing and building a circuits that generate audible
signals also expected to keep students interested, engaging and
motivating practical learning experience. Together, the three circuit
areas teach useful concepts in operational amplifier circuits, digital
logic, and filter design.

\section{Objective}\label{objective}

The primary objective is to develop synthesizer circuit architectures
that can be implemented on the LABOR platform, encompassing core
building blocks such as voltage-controlled oscillators (VCOs), envelope
generators, filters, and amplifiers.

The list of tasks to be accomplished in this thesis includes:

\begin{itemize}
\item
  To understand the electronic principles behind waveform conversion,
  sequencing and noise generation
\item
  To design and implement the circuits on the LABOR workstation
\item
  To test and evaluate the circuits, comparing measured results with
  theoretical expectations
\item
  To develop hardware /software interfaces for testing and analysing the
  circuits
\item
  To document the design process, circuit schematics, and experimental
  results in a clear and educational manner
\end{itemize}

\section{Work Packages}\label{work-packages}

\begin{enumerate}
\def\labelenumi{\arabic{enumi}.}
\item
  Technical research and literature review:

  This phase will focus on building a solid understanding of the
  technical topics like waveform conversion, sequencing, and noise
  generation. It will involve reviewing existing literature, tutorials,
  and research papers to gather relevant information and insights and to
  identify best practices and design principles for synthesizer circuit
  design.
\item
  Circuit design and implementation:

  In this phase, the knowledge gained from the research will be applied
  to design the synthesizer circuits to implement on the LABOR
  workstation. This will involve creating circuit schematics, selecting
  appropriate components. The design process will be iterative, with
  refinement of the circuits based on observed performance and
  theoretical expectations.
\item
  Testing and evaluation:

  This phase will involve testing the implemented circuits on the LABOR
  workstation, measuring their performance, and comparing the results
  with theoretical expectations. The testing process will include using
  oscilloscopes and other measurement tools to observe the generated
  signals. Each circuit section will first be examined independently so
  that possible problems can be identified more easily and depending on
  Depending on the circuit, data will be measured which include waveform
  shape, signal amplitude, timing accuracy, frequency response, and
  noise behaviour.
\item
  Software integration:

  In this phase, software tools will be developed to facilitate the
  testing and analysis of the circuits which include programming and
  designing the circuit using specific softare tools. The software will
  be used to visualize the generated signals, perform signal analysis.
\item
  Documentation:

  The final phase will bring together the results from the previous work
  packages and involve documenting the entire design process, circuit
  schematics, images, experimental results, and analysis in a clear and
  educational manner. The thesis will describe the development process
  and discuss both the successful results and the limitations
  encountered during the whole process.
\end{enumerate}

\section{Work Plan}\label{work-plan}

Use a Gantt-Chart with Mermaid, see: -
\href{https://quarto.org/docs/authoring/diagrams.html}{Quarto -
Diagrams} - \href{https://mermaid.js.org/syntax/gantt.html}{Mermaid -
Gantt}

Month 1: Technical research, literature review and analysing Moritz
Klein's tutorials. Month 2: Familiarization with the LABOR workstation,
initial circuit design, and circuit calculations. Month 3: Circuit
implementation on LABOR and testing of the individual circuits. Month 4:
Further testing, evaluation, and refinement of the circuits based on
observed performance. Month 5: Development of software tools for testing
and analysis, and integration with the hardware circuits Month 6:
Documentation of the entire design process, circuit schematics,
experimental results, and analysis.

\begin{figure}[H]

\centering{

\includegraphics[width=7.79in,height=3.79in]{chapter1_files/figure-latex/mermaid-figure-1.png}

}

\caption{\label{fig-gantt}Project timeline of the thesis work}

\end{figure}%

\section{References}\label{references}

Put them all in your bibliography (BibTeX file
\_bibliography/references.bib), use a reference manager like zotero and
then cite. This will automatically generate a biblikography at the end
of your thesis, see references.qmd file.

\%\%See: -
\href{https://quarto.org/docs/authoring/citations.html}{Quarto -
Citations} - \href{https://bibtex.eu/quarto/}{Quarto and BibTeX} -
\href{https://www.zotero.org}{Zotero} and
\href{https://github.com/retorquere/zotero-better-bibtex}{Better BibTeX}

\begin{enumerate}
\def\labelenumi{\arabic{enumi}.}
\tightlist
\item
  Modelling the Excess Noise due to Avalanche Multiplication in
  (Hetero-Junction) Bipolar Transistors - J. C. J. Paasschens and R. de
  Kort - Bipolar/BiCMOS Circuits and Technology, 2004.
\item
  Probability, Random Variables and Stochastic Processes - A. Papoulis,
  S. U. Pillai - 4th ed., McGraw-Hill, 2002.
\item
  The Differential Pair as a Triangle-Sine Wave Converter - R. G. Meyer,
  W. M. C. Sansen, S. Lui, S. Peeters - IEEE Journal of Solid-State
  Circuits, June 1976.
\item
  Design with Operational Amplifiers and Analog Integrated Circuits - S.
  Franco - 4th ed., McGraw-Hill, 2015
\item
  Limitations of a relaxation oscillator in capacitance measurements -
  Y. Liu, S. Chen, K. Nakayama, K. Watanabe - IEEE Transactions on
  Instrumentation and Measurement ( Volume: 49, Issue: 5), October 2000.
\item
  Tunable Waveform Generator for Sensor and Instrumentation Applications
  - Sougata Kumar Kar, Siddhartha Sen - IEEE Transactions on
  Instrumentation and Measurement ( Volume: 60, Issue: 10), October
  2011.
\item
  https://www.moritzkleininstruments.com/
\end{enumerate}

\emph{References are not only a list at the end of your report, they
need to be linked/referenced from your writings.}

\bookmarksetup{startatroot}

\chapter{Method}\label{sec-method}

This report underscores the transformative impact of reproducible
research on the scientific landscape, promoting a commitment to openness
and accountability for the benefit of the entire research community.

\bookmarksetup{startatroot}

\chapter*{References}\label{references-1}
\addcontentsline{toc}{chapter}{References}

\markboth{References}{References}

\protect\phantomsection\label{refs}
\begin{CSLReferences}{0}{0}
\bibitem[\citeproctext]{ref-Klein_2026}
\CSLLeftMargin{{[}1{]} }%
\CSLRightInline{M. Klein, {``LABOR.''} {[}Online{]}. Available:
\url{https://www.moritzkleininstruments.com}}

\bibitem[\citeproctext]{ref-Kar_2011}
\CSLLeftMargin{{[}2{]} }%
\CSLRightInline{S. K. Kar and S. Sen, {``Tunable square-wave generator
for integrated sensor applications,''} \emph{IEEE Transactions on
Instrumentation and Measurement}, vol. 60, no. 10, pp. 3369--3375, Oct.
2011, doi:
\href{https://doi.org/10.1109/tim.2011.2128490}{10.1109/tim.2011.2128490}.}

\bibitem[\citeproctext]{ref-Liu_2000}
\CSLLeftMargin{{[}3{]} }%
\CSLRightInline{Y. Liu, S. Chen, K. Nakayama, and K. Watanabe,
{``Limitations of a relaxation oscillator in capacitance
measurements,''} \emph{IEEE Transactions on Instrumentation and
Measurement}, vol. 49, no. 5, pp. 980--983, Oct. 2000, doi:
\href{https://doi.org/10.1109/19.872917}{10.1109/19.872917}.}

\bibitem[\citeproctext]{ref-Papoulis_2002}
\CSLLeftMargin{{[}4{]} }%
\CSLRightInline{A. Papoulis and S. U. Pillai, \emph{Probability, random
variables and stochastic processes}, 4th ed. McGraw-Hill, 2002, p. 186.}

\bibitem[\citeproctext]{ref-Franco_2001}
\CSLLeftMargin{{[}5{]} }%
\CSLRightInline{Sergio. Franco, \emph{Design with operational amplifiers
and analog integrated circuits}, 3rd ed. in McGraw-hill series in
electrical and computer engineering. McGraw-Hill, Inc., United States,
2001, p. 672. Available:
\url{https://dl.acm.org/doi/abs/10.5555/1593999}}

\bibitem[\citeproctext]{ref-Paasschens_2004}
\CSLLeftMargin{{[}6{]} }%
\CSLRightInline{J. C. J. Paasschens and R. de Kort, {``Modelling the
excess noise due to avalanche multiplication in (hetero-junction)
bipolar transistors,''} Institute of Electrical; Electronics Engineers
(IEEE), 2004, pp. 108--111. doi:
\href{https://doi.org/10.1109/BIPOL.2004.1365757}{10.1109/BIPOL.2004.1365757}.}

\end{CSLReferences}

\cleardoublepage
\phantomsection
\addcontentsline{toc}{part}{Appendices}
\appendix

\chapter{Tools}\label{sec-tools}

This thesis was written using \href{https://quarto.org/}{Quarto}

\chapter{USB-Stick and Script
Referencing}\label{usb-stick-and-script-referencing}

All directory structures and file paths presented in this appendix are
intentionally simplified. They are provided to clarify the organization
of the project and to indicate where specific scripts, netlists,
verification outputs, and layout files can be found within the design
environment.

\begin{Shaded}
\begin{Highlighting}[]
\NormalTok{thesis}
\NormalTok{├── .gitattributes}
\NormalTok{├── .gitignore}
\NormalTok{├── \_book/}
\NormalTok{├── \_quarto/}
\NormalTok{├── \_quarto.yml}
\NormalTok{├── \_bibliography}
\NormalTok{├── index.qmd            \# Main project page}
\NormalTok{├── README.md}
\NormalTok{├── references.qmd}
\NormalTok{├── appendix             \# Appendix}
\end{Highlighting}
\end{Shaded}



\backmatter


\end{document}
